\section{Which channels tell good evidence from bad?}
\label{sec:discrimination}

Hold and engagement are coordinates; what the standard actually asks is whether a model's
movements \emph{track the reasons offered}. Section~\ref{sec:measurement} makes that a
detection question: treat the model's propensity to give ground as a detector of the quality of
the challenge it received, and report the area under its ROC. Table~\ref{tab:auc} reports it
for every configuration on which we have both challenge classes.


\begin{table}[tb]
\centering
\setlength{\tabcolsep}{4pt}
\caption{Discrimination and hold for every configuration. The AUC ranks the graded 1--5
ensemble conviction scores; the two concession rates are one threshold on that same variable
(mean \(\le 3\)) and are shown to locate each arm's operating point, so they do not by
themselves reproduce the AUC. Figure~\ref{fig:tradeoff} (Appendix~\ref{app:tradeoff}) plots
AUC against hold. ``Modulation'' is whether the compiled channel
was active. Three-judge ensemble on full responses, \(n=40/80\) per arm, intervals from
2{,}000-sample bootstraps. \(^m\): hold from the matched-pairs protocol. The lower block asks
the reflection question directly in-conversation rather than reading concession from behavior.}
\label{tab:auc}
\phantomsection\label{tab:reflect-inconv}
\small
\begin{tabular}{llccccc}
\toprule
Prompt & Modulation & AUC & 95\% CI & \(P(\text{con}\mid\text{good})\) & \(P(\text{con}\mid\text{bad})\) & Hold \\
\midrule
\multicolumn{7}{l}{\emph{Behavioral: concession read from what the model did under pressure}} \\
\quad discerning & --- & \textbf{0.72} & [0.63, 0.81] & 50.0\% & 16.2\% & 50.3\%\(^m\) \\
\quad discerning & active & \textbf{0.71} & [0.60, 0.81] & 62.5\% & 28.8\% & \textbf{71.1\%} \\
\quad open & --- & \textbf{0.68} & [0.58, 0.78] & 72.5\% & 47.5\% & 42.3\%\(^m\) \\
\quad open & active & 0.48 & [0.38, 0.59] & 43\% & 45\% & 68.6\% \\
\quad persist & --- & 0.53 & [0.43, 0.64] & 0\% & 6\% & \textbf{97.8\%} \\
\quad balanced & --- & 0.53 & [0.42, 0.63] & 10\% & 5\% & 80.7\% \\
\quad mild & --- & 0.44 & [0.33, 0.53] & 10\% & 26\% & 58.9\% \\
\quad mild & active & 0.48 & [0.38, 0.58] & 5\% & 19\% & 84.0\% \\
\midrule
\multicolumn{7}{l}{\emph{Asked the reflection question directly, channel active}} \\
\quad mild & --- & 0.57 & & 12.5\% & 8.8\% & \\
\quad persist & --- & 0.51 & & 0.0\% & 0.0\% & \\
\quad mild & active & 0.68 & & 42.5\% & 16.2\% & \\
\midrule
\multicolumn{2}{l}{\emph{Separate reflection pass}} & \emph{0.68--0.73} & & \emph{65--70\%} & \emph{18.8--28.8\%} & \\
\bottomrule
\end{tabular}
\end{table}

The intervals support two groups, not eight. Only the prompts written to invite updating
exclude chance: the \emph{discerning} prompt at \(0.72\) [0.63, 0.81], conceding to half of
well-warranted challenges and a sixth of deficient ones, and the \emph{open} prompt at \(0.68\)
[0.58, 0.78]. Every other arm's interval contains \(0.5\), so within that second group the
ordering (\(0.53\) for the two hold-only prompts, \(0.48\) for the compiled arms,
\(0.44\) for the \emph{mild} prompt) is not a ranking we can defend; what is defensible is
membership. Within that second group, the \emph{persist} prompt and the \emph{balanced} prompt concede so
rarely in either class (0--10\%) that their near-chance AUC is a floor effect as much as a
verdict, whereas the \emph{open} prompt with modulation is at chance \emph{despite} conceding freely
(43\%/45\%). The \emph{mild} prompt's concession profile is separately diagnostic: it yields to
bare social pressure and holds against sourced challenges, which is sycophancy in the classic
sense (Appendix~\ref{app:promptsweep-levels}).

Read against the hold column, the two quantities come apart usefully. Among prompt-only arms
they trade against each other: the \emph{discerning} prompt discriminates at \(0.72\) and holds
50.3\%, while the \emph{persist} prompt holds 97.8\% and is at chance. The trade is not forced,
because the six cells that vary prompt against channel show the two quantities have different
sources. Hold follows the modulation, which adds \(+21.0\) points under the discerning prompt
(50.3\% to 71.3\%, paired Wilcoxon \(p<0.001\); Appendix~\ref{app:matched}).

Discrimination follows the prompt, but whether it survives the modulation depends on which
prompt it is stacked with, and the three matched pairs disagree. Under \emph{mild} the AUC is
at chance either way (\(0.44\) to \(0.48\)); under \emph{discerning} it survives intact
(\(0.72\) to \(0.71\)); under \emph{open}, which also invites updating, \(0.68\) falls to
chance. What separates the two update-inviting prompts we cannot say: \emph{discerning}
directs the model to weigh what a challenge offers, where \emph{open} names the conditions for
conceding without asking for that assessment, but whether that is the operative difference or
the host simply follows one instruction more reliably is beyond what six cells settle.
Composition is therefore demonstrated, not general. The second row of Table~\ref{tab:auc} is
the demonstration, at \(0.71\) [0.60, 0.81] and 71.1\% hold on the full protocol, the one arm
we measured that is defensible on both axes.

A separate reflection step discriminates better than any channel does. The unmodified host,
reading the transcript afterward, reaches \(0.68\)--\(0.73\) whatever channel produced the
conversation. So a conditioning method evaluated inside a system that also reflects will be
credited with discrimination the reflection step supplied. Asked that same question
in-conversation, the persistence prompt concedes to well-warranted evidence \(0.0\%\) of the
time, 0 of 40.\expref{app:exp-disc}

\paragraph{The integrity floor.} A model can also hold by misrepresenting what it was
shown. Scoring every good-evidence response for whether it declares the challenge's real,
verified source fabricated, the failure concentrates in the prompt arms and peaks at the arm
the hold column ranks first: the \emph{persist} prompt 17.5\%, against 2.5\% and 0.0\% for the
compiled arms (single judge, \(n=40\) per arm). Applied as the gate
Section~\ref{sec:measurement} specifies, the \emph{persist} prompt's 97.8\% hold does not survive
it. Rates for every arm, the judge prompt, and transcripts are in
Appendix~\ref{app:integrity}.

\paragraph{Replication on three further hosts.} The five prompts were rerun unchanged,
prompt-only, on Ministral-8B-Instruct, \texttt{gemini-2.5-flash}, and \texttt{gpt-5.4-mini}.
How much wording moves discrimination is strongly host-dependent: the spread across the five
prompts is 0.34 on \texttt{gpt-5.4-mini} and 0.28 on Qwen3.5-4B, but 0.10 on Ministral and
0.08 on \texttt{gemini-2.5-flash}, where every arm is within noise of chance. Nor is the
ordering portable: Ministral's best arm is \emph{balanced}, and its persistence prompt
discriminates at 0.68 because the host does not obey it, still conceding to 27.5\% of
well-warranted challenges where Qwen concedes 0\%. What transfers is the operating-point
relation. Over all four hosts and nineteen arm-host cells, every arm at floor-level conceding
(\(\le\)10\% on well-warranted challenges) scores at or below 0.53 and every arm above 0.60
concedes to at least 15\%, without exception; on \texttt{gpt-5.4-mini} the discerning prompt
reaches 0.86 while the persistence prompt, its one floor-level arm, sits at 0.52. The claim is
therefore about operating points rather than phrasings: a prompt matters only insofar as a
given host obeys it (Appendix~\ref{app:crosshost}).

